\documentclass[11pt,a4paper]{article}
\usepackage[T1]{fontenc}
\usepackage[utf8]{inputenc}
\usepackage[margin=2.5cm]{geometry}
\usepackage{tabularx}

\title{AforoSeguro\\Informe de práctica}
\author{Zambrano Parraga Luis Mario}
\date{23/11/2026}

\begin{document}
\maketitle

\section{Escenarios ejecutados}
Java 21. Un coordinador y tres puertas en la misma computadora, con aforo 10.
Las referencias de la tabla corresponden a transcripciones de consola.
E0 corresponde al control de una sesión supervisada y quedó fuera de este ensayo local.

\small
\begin{tabularx}{\textwidth}{@{} l l X @{}}
\textbf{Escenario} & \textbf{Evidencia} & \textbf{Observación} \\
E1 & \texttt{E1.txt} & Entró una persona por cada puerta y salió una por P2. Las tres consultas devolvieron 2/10. \\
E2 & \texttt{E2.txt} & La undécima entrada fue rechazada por aforo completo. El token incorrecto también fue rechazado y quedó registrado. \\
E3 & \texttt{E3.txt} & En modo normal se permitieron 10 entradas, se rechazaron 20 y el contador quedó en 10. En modo inseguro se permitieron 30, se rechazaron 0 y el contador quedó en 30. Hubo cero errores de comunicación. \\
E4 & \texttt{E4.txt} & Al detener P3, el coordinador la marcó como caída. Al reiniciarla, volvió a marcarla como activa. \\
E5 & \texttt{E5.txt} & Con el coordinador detenido, P1 indicó que la entrada quedó sin confirmar y pidió mantener el paso cerrado. \\
E6 & \texttt{E6.txt} & Se combinaron 432 eventos en el registro ordenado. \\
\end{tabularx}
\normalsize

\section{Causalidad}
En \texttt{logs/orden\_total.log} se pueden seguir dos solicitudes por su UUID.
Para el latido de P3 con prefijo \texttt{13159384}, P3 registró el envío
con marca 1 y el coordinador registró la recepción con marca 2.
Para la entrada de P1 con prefijo \texttt{0e830852}, P1 registró el envío
con marca 12 y el coordinador registró la recepción con marca 13.

\newpage

\section{Preguntas de análisis}
\begin{enumerate}
\item \textbf{Concurrencia.} En E3, varios hilos comprobaron el aforo antes
de que venciera la pausa de 50 ms. Cada uno pudo incrementar el contador
después de haber observado espacio libre. En modo normal, la exclusión
mutua abarca tanto la comprobación como la actualización. Los hilos
pueden atender conexiones en paralelo y esperan su turno para modificar
el contador.

\item \textbf{Coordinador inaccesible.} En E5, P1 mantuvo el paso cerrado.
Ante una interrupción de la comunicación, la puerta sacrifica la
disponibilidad de las entradas para preservar la consistencia del aforo,
que es el aspecto relevante de CAP en este escenario. El público debe
esperar y la administración debe restablecer la coordinación antes de
autorizar el acceso.

\item \textbf{Punto único de fallo.} Una elección como Bully puede designar
un nuevo coordinador. El sucesor necesita un contador replicado y un
mecanismo para distinguir las operaciones ya aplicadas. Dos líderes
activos a la vez podrían aceptar entradas a partir de estados
diferentes. En este ensayo, reiniciar el coordinador puso el contador
en cero, aunque las bitácoras quedaron en disco.

\item \textbf{Latencia y latidos.} El intervalo de 2 segundos y el umbral
de 6 segundos permiten detectar la ausencia de varios latidos. Esto
toma en cuenta la falacia de que la latencia es cero. Con 7 segundos
de latencia, el coordinador podría marcar como caída una puerta sana.
El cliente también tiene un límite de espera de 3 segundos para la
respuesta, por lo que una solicitud podría quedar sin confirmar incluso
si el coordinador la aplica más tarde.
\end{enumerate}

\end{document}
